% from ega1-10-fr.tex
\oldpage[I]{181}
l'idéal de définition $\mathfrak{J}$ considéré~; notons-le
$\mathfrak{X}$. Soit $(\mathfrak{J}_\lambda)$ un système fondamental de
voisinages de 0 dans $A$, formé d'idéaux de définition, et pour tout
$\lambda$, soit $\mathscr{O}_\lambda$ le faisceau structural de
$\operatorname{Spec}(A/\mathfrak{J}_\lambda)$~; ce faisceau est induit sur
$\mathfrak{X}$ par
$\widetilde{A}/\widetilde{\mathfrak{J}}_\lambda$ (lequel est nul hors de
$\mathfrak{X}$).

Pour $\mathfrak{J}_\mu\subset\mathfrak{J}_\lambda$, l'homomorphisme
canonique $A/\mathfrak{J}_\mu\to A/\mathfrak{J}_\lambda$ définit donc un
homomorphisme
$u_{\lambda\mu}:\mathscr{O}_\mu\to\mathscr{O}_\lambda$ de faisceaux
d'anneaux (\hyperref[I.1.6.1-fr]{1.6.1}), et
$(\mathscr{O}_\lambda)$ est un \emph{système projectif de faisceaux
d'anneaux} pour ces homomorphismes. Comme la topologie de $\mathfrak{X}$
admet une base formée d'ouverts quasi-compacts, on peut associer à tout
$\mathscr{O}_\lambda$ un \emph{faisceau d'anneaux topologiques
pseudo-discrets} (\hyperref[0.3.8.1-fr]{0, 3.8.1}) qui a
$\mathscr{O}_\lambda$ comme faisceau d'anneaux (sans topologie) sous-jacent,
et que nous noterons encore $\mathscr{O}_\lambda$~; et les
$\mathscr{O}_\lambda$ forment encore un \emph{système projectif de faisceaux
d'anneaux topologiques} (\hyperref[0.3.8.2-fr]{0, 3.8.2}). Nous désignerons
par $\mathscr{O}_{\mathfrak{X}}$ le \emph{faisceau d'anneaux topologiques}
sur $\mathfrak{X}$, limite projective du système
$(\mathscr{O}_\lambda)$~; pour tout ouvert \emph{quasi-compact} $U$ de
$\mathfrak{X}$, $\Gamma(U,\mathscr{O}_{\mathfrak{X}})$ est donc l'anneau
topologique limite projective du système d'anneaux \emph{discrets}
$\Gamma(U,\mathscr{O}_\lambda)$
(\hyperref[0.3.2.6-fr]{0, 3.2.6}).
\end{env}

\begin{definition}[10.1.2]
\label{I.10.1.2-fr}
Étant donné un anneau topologique admissible $A$, on appelle spectre formel
de $A$ et on note $\operatorname{Spf}(A)$ le sous-espace fermé
$\mathfrak{X}$ de $\operatorname{Spec}(A)$ formé des idéaux premiers
ouverts de $A$. On dit qu'un espace topologiquement annelé est un schéma
formel affine s'il est isomorphe à un spectre formel
$\operatorname{Spf}(A)=\mathfrak{X}$ muni du faisceau d'anneaux
topologiques $\mathscr{O}_{\mathfrak{X}}$ limite projective des faisceaux
d'anneaux pseudo-discrets
$(\widetilde{A}/\widetilde{\mathfrak{J}}_\lambda)|\mathfrak{X}$, où
$\mathfrak{J}_\lambda$ parcourt l'ensemble filtrant des idéaux de
définition de $A$.
\end{definition}

Quand nous parlerons d'un \emph{spectre formel}
$\mathfrak{X}=\operatorname{Spf}(A)$ comme d'un schéma formel affine, il
s'agira toujours de l'espace topologiquement annelé
$(\mathfrak{X},\mathscr{O}_{\mathfrak{X}})$ où
$\mathscr{O}_{\mathfrak{X}}$ est défini comme ci-dessus.

On notera que tout \emph{schéma affine} $X=\operatorname{Spec}(A)$ peut être
considéré comme un schéma formel affine d'une seule manière, en considérant
$A$ comme un anneau topologique discret~: les anneaux topologiques
$\Gamma(U,\mathscr{O}_X)$ sont alors discrets lorsque $U$ est
quasi-compact (mais non en général lorsque $U$ est un ouvert quelconque de
$X$).

\begin{proposition}[10.1.3]
\label{I.10.1.3-fr}
Si $\mathfrak{X}=\operatorname{Spf}(A)$, où $A$ est un anneau admissible,
$\Gamma(\mathfrak{X},\mathscr{O}_{\mathfrak{X}})$ est topologiquement
isomorphe à $A$.
\end{proposition}

\begin{proof}
En effet, comme $\mathfrak{X}$ est fermé dans $\operatorname{Spec}(A)$, il
est quasi-compact, et par suite
$\Gamma(\mathfrak{X},\mathscr{O}_{\mathfrak{X}})$ est topologiquement
isomorphe à la limite projective des anneaux discrets
$\Gamma(\mathfrak{X},\mathscr{O}_\lambda)$~; mais
$\Gamma(\mathfrak{X},\mathscr{O}_\lambda)$ est isomorphe à
$A/\mathfrak{J}_\lambda$ (\hyperref[I.1.3.7-fr]{1.3.7})~; comme $A$ est
séparé et complet, il est topologiquement isomorphe à
$\varprojlim A/\mathfrak{J}_\lambda$
(\hyperref[0.7.2.1-fr]{0, 7.2.1}), d'où la proposition.
\end{proof}

\begin{proposition}[10.1.4]
\label{I.10.1.4-fr}
Soient $A$ un anneau admissible, $\mathfrak{X}=\operatorname{Spf}(A)$, et
pour tout $f\in A$, soit
$\mathfrak{D}(f)=D(f)\cap\mathfrak{X}$~; l'espace topologiquement annelé
$(\mathfrak{D}(f),\mathscr{O}_{\mathfrak{X}}|\mathfrak{D}(f))$ est
isomorphe au spectre formel affine
$\operatorname{Spf}(A_{\{f\}})$
(\hyperref[0.7.6.15-fr]{0, 7.6.15}).
\end{proposition}

\begin{proof}
Pour tout idéal de définition $\mathfrak{J}$ de $A$, l'anneau discret
$S_f^{-1}A/S_f^{-1}\mathfrak{J}$ s'identifie canoniquement à
$A_{\{f\}}/\mathfrak{J}_{\{f\}}$
(\hyperref[0.7.6.9-fr]{0, 7.6.9}), donc
(\hyperref[I.1.2.5-fr]{1.2.5} et \hyperref[I.1.2.6-fr]{1.2.6}), l'espace
topologique $\operatorname{Spf}(A_{\{f\}})$ s'identifie canoniquement à
$\mathfrak{D}(f)$. En outre, pour tout ouvert quasi-compact $U$ de
$\mathfrak{X}$ contenu dans $\mathfrak{D}(f)$,
$\Gamma(U,\mathscr{O}_\lambda)$ s'identifie au module des sections du
faisceau structural de
$\operatorname{Spec}(S_f^{-1}A/S_f^{-1}\mathfrak{J}_\lambda)$ au-dessus de
$U$ (\hyperref[I.1.3.6-fr]{1.3.6}), donc, si on pose
$\mathfrak{Y}=\operatorname{Spf}(A_{\{f\}})$,
$\Gamma(U,\mathscr{O}_{\mathfrak{X}})$ s'identifie au module des sections
$\Gamma(U,\mathscr{O}_{\mathfrak{Y}})$, ce qui démontre la proposition.
\end{proof}

